\documentclass[11pt,aspectratio=169]{ctexbeamer}

\usetheme{default}
\setbeamertemplate{navigation symbols}{}

\usepackage{amsmath}
\usepackage{booktabs}

% Projector-safe palette: high contrast on common lecture-room projectors.
\definecolor{DeepBlue}{RGB}{11,58,102}
\definecolor{SoftInk}{RGB}{31,41,51}
\definecolor{PaperBg}{RGB}{255,255,255}
\definecolor{PanelBlue}{RGB}{226,237,247}

\setbeamercolor{normal text}{fg=SoftInk,bg=PaperBg}
\setbeamercolor{structure}{fg=DeepBlue}
\setbeamercolor{frametitle}{fg=DeepBlue,bg=PaperBg}
\setbeamercolor{title}{fg=DeepBlue}
\setbeamercolor{block title}{fg=DeepBlue,bg=PanelBlue}
\setbeamercolor{block body}{fg=SoftInk,bg=white}

\newcommand{\Core}[1]{
  \begin{block}{核心结论}
    #1
  \end{block}
}

\newcommand{\LectureMeta}[2]{
  \title{大模型推理系统与实践}
  \subtitle{#1}
  \author{张书豪}
  \institute{华中科技大学计算机学院}
  \date{2027 · #2}
}

\newcommand{\PaperRef}[3]{
  \item \textbf{#1}（#2）\hfill\href{#3}{原文链接}
}

\newcommand{\ResearchSources}[1]{
  \begin{frame}{本讲参考资料}
    \normalsize
    \begin{itemize}
      #1
    \end{itemize}
    \vspace{0.3em}
    \textcolor{SoftInk}{课件中的论文数字均指作者在特定模型、硬件和工作负载下报告的结果，不代表无条件性能。}
  \end{frame}
}

\newcommand{\InternalCaseNote}{
  \small\textcolor{SoftInk}{说明：本页为课程研究案例或教学假设，不代表已发表论文结论；只有在实验包给出模型、硬件、工作负载和原始记录时，数字才可作为结果引用。}
}

% 14 讲细分版的统一参考资料入口。论文数字的解释边界见课程内容基线。

\newcommand{\LectureOneSources}{\ResearchSources{
  \PaperRef{vLLM: Easy, Fast, and Cheap LLM Serving with PagedAttention}{SOSP 2023}{https://doi.org/10.1145/3600006.3613165}
  \PaperRef{SGLang: Efficient Execution of Structured Language Model Programs}{arXiv:2312.07104}{https://arxiv.org/abs/2312.07104}
  \PaperRef{nano-vllm-hust 教学代码}{课程仓库}{https://github.com/intellistream/intro-to-llm-inference-engines}
}}

\newcommand{\LectureTwoSources}{\ResearchSources{
  \PaperRef{vLLM / PagedAttention}{SOSP 2023}{https://doi.org/10.1145/3600006.3613165}
  \PaperRef{Sarathi-Serve}{OSDI 2024}{https://www.usenix.org/conference/osdi24/presentation/agrawal}
  \PaperRef{DistServe}{OSDI 2024}{https://www.usenix.org/conference/osdi24/presentation/zhong-yinmin}
}}

\newcommand{\LectureThreeSources}{\ResearchSources{
  \PaperRef{vLLM / PagedAttention}{SOSP 2023}{https://doi.org/10.1145/3600006.3613165}
  \PaperRef{DistServe}{OSDI 2024}{https://www.usenix.org/conference/osdi24/presentation/zhong-yinmin}
  \PaperRef{Llumnix}{OSDI 2024}{https://www.usenix.org/conference/osdi24/presentation/sun-biao}
}}

\newcommand{\LectureFourSources}{\ResearchSources{
  \PaperRef{Orca}{OSDI 2022}{https://www.usenix.org/conference/osdi22/presentation/yu}
  \PaperRef{Sarathi-Serve}{OSDI 2024}{https://www.usenix.org/conference/osdi24/presentation/agrawal}
  \PaperRef{FastServe}{NSDI 2026}{https://www.usenix.org/conference/nsdi26/presentation/wu-bingyang}
}}

\newcommand{\LectureFiveSources}{\ResearchSources{
  \PaperRef{vLLM / PagedAttention}{SOSP 2023}{https://doi.org/10.1145/3600006.3613165}
  \PaperRef{vAttention}{arXiv:2405.04437}{https://arxiv.org/abs/2405.04437}
  \PaperRef{Rethinking KV Cache Compression}{MLSys 2025}{https://proceedings.mlsys.org/paper_files/paper/2025/hash/26289c647c6828e862e271ca3c490486-Abstract-Conference.html}
  \PaperRef{Mooncake}{FAST 2025}{https://www.usenix.org/conference/fast25/presentation/qin}
}}

\newcommand{\LectureSixSources}{\ResearchSources{
  \PaperRef{DistServe}{OSDI 2024}{https://www.usenix.org/conference/osdi24/presentation/zhong-yinmin}
  \PaperRef{RAGCache}{arXiv:2404.12457}{https://arxiv.org/abs/2404.12457}
  \PaperRef{Mooncake}{FAST 2025}{https://www.usenix.org/conference/fast25/presentation/qin}
  \PaperRef{NEO}{MLSys 2025}{https://proceedings.mlsys.org/paper_files/paper/2025/hash/66a026c0d17040889b50f0dfa650e5e0-Abstract-Conference.html}
}}

\newcommand{\LectureSevenSources}{\ResearchSources{
  \PaperRef{vLLM / PagedAttention}{SOSP 2023}{https://doi.org/10.1145/3600006.3613165}
  \PaperRef{SGLang}{arXiv:2312.07104}{https://arxiv.org/abs/2312.07104}
  \PaperRef{DistServe}{OSDI 2024}{https://www.usenix.org/conference/osdi24/presentation/zhong-yinmin}
  \PaperRef{Mooncake}{FAST 2025}{https://www.usenix.org/conference/fast25/presentation/qin}
}}

\newcommand{\LectureEightSources}{\ResearchSources{
  \PaperRef{FlashAttention}{NeurIPS 2022}{https://arxiv.org/abs/2205.14135}
  \PaperRef{FlexGen}{ICML 2023}{https://proceedings.mlr.press/v202/sheng23a.html}
  \PaperRef{PowerInfer}{SOSP 2024 / arXiv:2312.12456}{https://arxiv.org/abs/2312.12456}
  \PaperRef{FlashInfer}{MLSys 2025}{https://proceedings.mlsys.org/paper_files/paper/2025/hash/dbf02b21d77409a2db30e56866a8ab3a-Abstract-Conference.html}
}}

\newcommand{\LectureNineSources}{\ResearchSources{
  \PaperRef{DEEPSERVE}{USENIX ATC 2025}{https://www.usenix.org/conference/atc25/presentation/hu-junhao}
  \PaperRef{QServe}{MLSys 2025}{https://proceedings.mlsys.org/paper_files/paper/2025/hash/fbe2b2f74a2ece8070d8fb073717bda6-Abstract-Conference.html}
  \PaperRef{vLLM-HUST}{开源维护仓库}{https://github.com/vLLM-HUST/vllm-hust}
}}

\newcommand{\LectureTenSources}{\ResearchSources{
  \PaperRef{vLLM / PagedAttention}{SOSP 2023}{https://doi.org/10.1145/3600006.3613165}
  \PaperRef{Revisiting Disaggregated LLM Serving}{2025 preprint, arXiv:2601.08833}{https://arxiv.org/abs/2601.08833}
  \PaperRef{Rethinking KV Cache Compression}{MLSys 2025}{https://proceedings.mlsys.org/paper_files/paper/2025/hash/26289c647c6828e862e271ca3c490486-Abstract-Conference.html}
}}

\newcommand{\LectureElevenSources}{\ResearchSources{
  \PaperRef{MLPerf Inference}{MLCommons benchmark}{https://mlcommons.org/benchmarks/inference-datacenter/}
  \PaperRef{ACM Artifact Review and Badging}{ACM policy}{https://www.acm.org/publications/policies/artifact-review-and-badging-current}
  \PaperRef{Revisiting Disaggregated LLM Serving}{2025 preprint, arXiv:2601.08833}{https://arxiv.org/abs/2601.08833}
}}

\newcommand{\LectureTwelveSources}{\ResearchSources{
  \PaperRef{nano-vllm-hust}{课程教学代码}{https://github.com/intellistream/intro-to-llm-inference-engines}
  \PaperRef{vLLM-HUST}{HUST 维护的 vLLM fork}{https://github.com/vLLM-HUST/vllm-hust}
  \PaperRef{上游 vLLM}{开源项目}{https://github.com/vllm-project/vllm}
}}

\newcommand{\LectureThirteenSources}{\ResearchSources{
  \PaperRef{课程内容基线与实验材料}{课程仓库}{https://github.com/intellistream/intro-to-llm-inference-engines}
  \PaperRef{DEEPSERVE}{USENIX ATC 2025}{https://www.usenix.org/conference/atc25/presentation/hu-junhao}
  \PaperRef{GhostServe}{MLSys 2026}{https://proceedings.mlsys.org/paper_files/paper/2026/hash/d99e8e80a6c41e148db686918dd7eab3-Abstract-Conference.html}
}}

\newcommand{\LectureFourteenSources}{\ResearchSources{
  \PaperRef{课程完整材料与后续更新}{课程仓库}{https://github.com/intellistream/intro-to-llm-inference-engines}
  \PaperRef{SpecGen}{MLSys 2026}{https://proceedings.mlsys.org/paper_files/paper/2026/hash/66a026c0d17040889b50f0dfa650e5e0-Abstract-Conference.html}
  \PaperRef{GhostServe}{MLSys 2026}{https://proceedings.mlsys.org/paper_files/paper/2026/hash/d99e8e80a6c41e148db686918dd7eab3-Abstract-Conference.html}
}}



\LectureMeta{第 2 讲 工作负载与评价指标}{研究生课程课件}

\begin{document}

\begin{frame}[plain]
  \titlepage
\end{frame}

\begin{frame}{教学目标}
  \Core{本讲围绕工作负载与评价指标展开，建立刻画推理系统问题的基本分析坐标。}
  \begin{itemize}
    \item 区分离线、在线、RAG 与智能体四类典型服务场景。
    \item 理解 TTFT、TPOT、P99 与 goodput 等核心指标的系统含义。
    \item 说明单点指标为何不足以代表端到端系统收益。
  \end{itemize}
\end{frame}

\begin{frame}{本讲的基本关系}
  \Core{第二讲的核心关系是：工作负载先定义问题边界，评价指标再定义优化目标与解释方式。}
  \begin{itemize}
    \item 相同指标在不同工作负载下可能对应不同的系统含义。
    \item 相同 workload 在不同优化目标下也可能导向不同的系统设计。
    \item 因而所有实验和比较都需要先交代 workload，再报告指标。
  \end{itemize}
\end{frame}

\begin{frame}{典型工作负载类型}
  \Core{不同工作负载会重新定义系统瓶颈，因此系统设计不能脱离场景。}
  \begin{itemize}
    \item 离线批推理关注总吞吐与资源利用率。
    \item 在线对话关注 TTFT、TPOT 与尾时延。
    \item RAG 服务关注检索到生成的全链路时延。
    \item 智能体工作流关注跨轮状态、工具调用与任务级完成时间。
  \end{itemize}
\end{frame}

\begin{frame}{四类工作负载的系统差异}
  \Core{同样被称为“推理”的服务，在不同 workload 下会呈现显著不同的主导矛盾。}
  \begin{itemize}
    \item 离线推理更偏向吞吐与成本优化。
    \item 在线对话更偏向 TTFT 与尾时延控制。
    \item RAG 更偏向检索与生成链路协同。
    \item 智能体更偏向多阶段流水与跨轮状态维护。
  \end{itemize}
\end{frame}

\begin{frame}{为什么工作负载变了，系统问题也会变}
  \Core{长上下文、MoE、多模态与智能体工作流会从容量、调度、通信与状态维护多个层面放大系统复杂度。}
  \begin{itemize}
    \item 长上下文放大 prefill 代价与 KV 压力。
    \item MoE 放大负载倾斜与通信波动。
    \item 智能体将单轮生成扩展为多阶段流水与持续写入问题。
  \end{itemize}
\end{frame}

\begin{frame}{长上下文如何改变瓶颈}
  \Core{长上下文会同时推高 prefill 计算代价、KV 驻留压力与排队放大效应。}
  \begin{itemize}
    \item prompt 变长会直接增加 prefill 开销。
    \item KV 驻留时间更长，状态占用更大。
    \item 长短请求混部时更容易出现尾时延恶化。
  \end{itemize}
\end{frame}

\begin{frame}{核心评价指标}
  \Core{评价指标的选择实质上对应系统优化目标的定义。}
  \begin{itemize}
    \item TTFT：首 token 等待时间。
    \item TPOT：相邻输出 token 的平均间隔。
    \item Throughput：单位时间完成的请求数或 token 数。
    \item P95/P99：尾时延表现。
    \item Goodput：满足服务目标的有效吞吐。
  \end{itemize}
\end{frame}

\begin{frame}{goodput 的形式化表达}
  \Core{goodput 的关键不在于重新命名 throughput，而在于把服务目标是否满足纳入统计口径。}
  \[
    \text{Goodput} = \frac{\#\{\text{requests meeting SLO}\}}{\Delta t}
  \]
  \begin{itemize}
    \item 当大量请求违反 SLO 时，原始 throughput 的解释价值会迅速下降。
    \item goodput 更适合比较调度、缓存与资源治理策略的真实服务效果。
  \end{itemize}
\end{frame}

\begin{frame}{指标之间为什么会冲突}
  \Core{指标冲突是系统优化中的常态，因此必须首先明确主目标与可接受代价。}
  \begin{itemize}
    \item 更大的 batch 通常提高吞吐，但可能恶化 TTFT。
    \item 更激进的缓存策略可能提高命中率，但也提高管理代价。
    \item 更高设备利用率并不等于更低 P99。
  \end{itemize}
\end{frame}

\begin{frame}{为什么平均值常常不够}
  \Core{用户体验和服务质量常由尾部与波动决定，而不是平均值。}
  \begin{itemize}
    \item 更大的 batch 可能提高平均吞吐，但也可能恶化 TTFT。
    \item 更高命中率未必自动带来端到端时延下降。
    \item 单次最好结果并不等于稳定收益。
  \end{itemize}
\end{frame}

\begin{frame}{为什么要引入 goodput}
  \Core{goodput 的意义在于把“服务目标是否满足”纳入吞吐定义，而不是只统计原始完成量。}
  \begin{itemize}
    \item 若大量请求违反 SLO，原始 throughput 的意义会迅速下降。
    \item goodput 更适合比较不同调度或资源治理策略。
    \item 它将系统目标从“完成更多”转化为“有效完成更多”。
  \end{itemize}
\end{frame}

\begin{frame}{工作负载与指标的先后关系}
  \Core{指标不能脱离场景单独解释；只有先交代 workload，数字才具有明确含义。}
  \begin{itemize}
    \item 离线批处理更接近资源利用率问题，在线对话更接近响应性问题。
    \item RAG 与 Agent 会把“生成”嵌入更长链路，因此单点吞吐不再足以代表体验。
    \item 同一套指标在不同 workload 下会具有不同含义。
  \end{itemize}
\end{frame}

\begin{frame}{例子：不同在线 workload 的差异}
  \Core{即便都属于在线请求，不同 workload 的长度分布、链路结构与服务目标也可能完全不同。}
  \small
  \begin{tabular}{p{0.20\textwidth}p{0.29\textwidth}p{0.31\textwidth}}
    \toprule
    场景 & 分布特征 & 更敏感的系统问题 \\
    \midrule
    对话服务 & 到达连续，prompt 长度波动大 & TTFT、尾时延、短请求体验 \\
    RAG 服务 & 检索与生成交织，链路更长 & 全链路时延、阶段协同 \\
    Agent 服务 & 多轮工具调用与多次生成 & 任务级完成时间、跨轮状态维护与一致性 \\
    \bottomrule
  \end{tabular}
\end{frame}

\begin{frame}{例子：均值为何经常产生误导}
  \Core{系统实验中最常见的误区，是只看均值而忽视分布、尾部与服务目标。}
  \begin{tabular}{p{0.26\textwidth}p{0.28\textwidth}p{0.28\textwidth}}
    \toprule
    表面现象 & 可能遗漏的信息 & 更合适的补充指标 \\
    \midrule
    平均 TTFT 正常 & 长请求与突发到达已推高 P99 & P95/P99、队列等待 \\
    平均吞吐更高 & 短请求体验可能显著恶化 & TTFT、短请求完成时间 \\
    完成量增加 & 满足 SLO 的请求比例可能下降 & Goodput、SLO 达成率 \\
    \bottomrule
  \end{tabular}
\end{frame}

\begin{frame}{代码例子：benchmark 如何固定分布并生成指标}
  \Core{工作负载与指标分析需要学生看到：实验结果并非“随便跑一遍”即可解释。}
  \begin{columns}[T,totalwidth=\textwidth]
    \column{0.5\textwidth}
    \begin{block}{输入分布控制}
      {\ttfamily\small
      \detokenize{seed(0)}\par
      \detokenize{prompt_token_ids = [[randint(...) ...]]}\par
      \detokenize{sampling_params = [SamplingParams(...)]}\par
      }
    \end{block}
    \column{0.46\textwidth}
    \begin{block}{结果计算逻辑}
      {\ttfamily\small
      \detokenize{t = time.time()}\par
      \detokenize{llm.generate(..., use_tqdm=False)}\par
      \detokenize{throughput = total_tokens / (time.time() - t)}\par
      }
    \end{block}
  \end{columns}
  \vspace{0.4em}
  \begin{itemize}
    \item 输入长度分布与输出长度分布需要分别说明。
    \item 仅有总体 throughput 仍不足以支撑 TTFT、P99 与队列等待层面的完整结论。
  \end{itemize}
\end{frame}

\begin{frame}{例子：throughput、TTFT 与 P99 的张力}
  \Core{工作负载与评价指标这一讲的重点，是接受真实系统中的指标张力并据此陈述取舍。}
  \begin{itemize}
    \item 更大的 batch 往往提高总体吞吐，但也可能抬高等待时间与 TTFT。
    \item 更保守的调度可能稳定尾时延，却压低设备利用率与平均吞吐。
    \item 因而任何优化都应首先说明优先保留的指标与可接受的代价。
  \end{itemize}
\end{frame}

\begin{frame}{课堂练习：为同一系统写出不同的指标解释}
  \Core{同一组数字可以对应不同解释，关键在于问题定义、工作负载设定与目标指标选择。}
  \begin{enumerate}
    \item 若一个系统吞吐提高 15\%，但 TTFT 恶化 30\%，如何判断它是否更适合在线服务。
    \item 给出一个适用于 Agent 工作流的 SLO 组合，而不只写单个 TTFT 阈值。
    \item 说明为什么指标解释必须回到 workload 与场景。
  \end{enumerate}
\end{frame}


\begin{frame}{课堂讨论}
  \Core{指标选择本身构成系统问题定义的重要组成部分。}
  \begin{itemize}
    \item 为什么以 TTFT 为核心的服务不会只优化峰值吞吐。
    \item RAG 服务与对话服务的主指标为何不同。
    \item 什么情况下应报告 goodput 而非 raw throughput。
  \end{itemize}
\end{frame}

\begin{frame}{课堂练习}
  \Core{本讲练习采用“场景 - 指标 - 瓶颈”三列表格式，以训练问题刻画能力。}
  \begin{itemize}
    \item 每个场景至少写出两个主指标。
    \item 每个指标至少写出一个可能冲突指标。
    \item 每个冲突都要对应一个系统机制取舍。
  \end{itemize}
\end{frame}

\begin{frame}{本讲总结}
  \Core{明确工作负载与评价指标，是后续讨论调度、KV 与执行优化问题的共同前提。}
  \begin{itemize}
    \item 不同场景对应不同瓶颈。
    \item 指标定义决定优化方向与论文叙事边界。
  \end{itemize}
\end{frame}

\LectureTwoSources

\begin{frame}{对应 Tutorial}
  \Core{本讲对应 Tutorial 1。先独立作答，再使用对应 Tutorial 的教师讨论要点核对。}
  \begin{itemize}
    \item 文件：\texttt{\detokenize{build/tutorials/Tutorial_01_工作负载与评价指标.pdf}}
    \item 动手运行、日志记录与提交要求统一查看 \texttt{\detokenize{build/experiments/}} 和对应 experiment 源目录下的 \texttt{\detokenize{assignment_spring-2027/}}。
  \end{itemize}
\end{frame}

\end{document}
